\section{Actual original tensor sections preserve true norms and nonzero sections}
\label{sec:scheme-global-section-tensor-product-norms}

\begin{theorem}[Every true original global-section tensor norm multiplies]
\label{thm:scheme-global-section-tensor-norm}
\lean{EnvelopeProperty.SchemeGlobalSectionTensorProductNorms.globalSectionNorm_tensorSection}
\leanok
\uses{thm:scheme-global-section-tensor-power-norm}
Let the actual original scheme be proper over the actual bounded valuation
base. For any two original invertible modules and any two actual entire global
sections, the true original proper-model norm of their actual tensor section
is the product of their true original norms, at every original native point.
Zero sections and zero norms are included.
\end{theorem}
\begin{proof}
The actual original tensor-section morphism is the tensor of the two original
structure-to-line morphisms after the inverse structure-unit multiplication.
The genuine original proper-model morphism norm multiplies on this tensor and
the structure-unit isomorphism has norm one. The actual whole-section tensor
identity identifies precisely the canonical original tensor global section.
\end{proof}

\begin{theorem}[Actual tensor-power sections retain both whole original factors]
\label{thm:scheme-global-section-power-tensor-norm}
\lean{EnvelopeProperty.SchemeGlobalSectionTensorProductNorms.powerTensorSection,EnvelopeProperty.SchemeGlobalSectionTensorProductNorms.globalSectionNorm_powerTensorSection}
\leanok
\uses{thm:scheme-global-section-tensor-norm,thm:scheme-global-section-model-norm-isometry,lem:fc-original-bundle-tensor-power-tensor}
For every original nonnegative integer $d$, two actual whole original sections
of $M^{\otimes d}$ and $N^{\otimes d}$ give an actual whole section of
$(M\otimes N)^{\otimes d}$. Its true norm at every original native point is
exactly the product of the original factor-section norms.
\end{theorem}
\begin{proof}
Take the actual original tensor section and carry it through the inverse
actual whole power-tensor isomorphism. This isomorphism is constructed from
true tensor associativity, symmetry and unit maps. The actual whole-section
isomorphism preserves the true proper-model norm. Apply the preceding
whole-section tensor norm identity.
\end{proof}

\begin{theorem}[Original nonzero tensor sections are derived intrinsically]
\label{thm:scheme-global-section-tensor-nonzero}
\lean{EnvelopeProperty.SchemeGlobalSectionTensorProductNorms.tensorSection_eq_zero_iff,EnvelopeProperty.SchemeGlobalSectionTensorProductNorms.powerTensorSection_eq_zero_iff}
\leanok
\uses{thm:scheme-global-section-tensor-norm,thm:scheme-global-section-power-tensor-norm,thm:scheme-original-global-section-norm-detection}
Assume that the actual original scheme $X$ is integral and proper over the
original field. The actual tensor section is zero if and only if one original
factor section is zero. The same exact criterion holds for the preceding
actual tensor-power section product. No nonzero-product or analytic
faithfulness certificate is assumed.
\end{theorem}
\begin{proof}
Both original nonzero sections have actual norm one at the same genuine
native point supported at the original generic point with trivial residue
valuation. The product norm is therefore one there, so the actual tensor
section cannot be zero. If a factor is zero, its actual norm vanishes at
every original native point; the product norm vanishes everywhere, and true
original norm detection forces the tensor section to be zero. The actual
power-tensor section norm identity gives the identical argument for powers.
These statements supply genuine original finite-family tensor calculus;
they do not assert independence of arbitrary numerical representatives or
finite dimensionality of the whole intrinsic numerical space.
\end{proof}
